\documentclass[english, inproceedings]{polze}

\begin{document}

\author{Tim Felgentreff}
\group{Software Architecture Group}
\organisation{Hasso-Plattner-Institut}
\email{tim.felgentreff@hpi.uni-potsdam.de}
\event{Fall 2013 Workshop}

\title{Solving Multidomain Constraints on Object Behavior}
\maketitle
\begin{abstract}
  Constraints allow developers to specify desired properties of a
  system and have those properties be maintained automatically.  This
  results in compact, localized code, avoiding scattered code to
  trigger checks and imperatively re-satisfy invariants.  Despite the
  improved comprehensibility and maintainability, as well as possible
  performance advantages, constraint programming is not yet
  widespread, with standard imperative programming still the norm.

  We propose an extension to existing imperative, object-oriented
  languages that unifies the methods of encapsulation and abstraction
  for both the declarative and the imperative paradigm. Our current
  research addresses integration of these different paradigms into one
  another and how this integration can be useful for in a variety of
  application domains.
\end{abstract}

\section{Introduction}
\label{sec:introduction}
Constraint programming has a long history of research interest, dating
back to the SketchPad~\cite{sutherland-ifips-1963} system. As a
paradigm for general purpose programming, the
CLP/R~\cite{jaffar-toplas-1992} and Kaleidoscope~\cite{lopez-phd}
systems provided integration of constraints with logic programming and
objects, respectively. Recently, there has been renewed interest
in constraints in programming, through libraries or in DSLs such as
the in the Mac OS X layout system.  Constraint programming allows
developers to precisely express desired properties of a system in a
number of domains, and have those properties be automatically
maintained by the runtime. This enables concise code, while avoiding
scattered code that trigger checks and re-satisfaction of those
properties.  (As an example, consider a draggable rectangle in a
graphical application that should follow the mouse when dragged, but
not leave the boundaries of its container. (continue this?))

Existing approaches to integrating constraint programming with the
imperative paradigm tried to bring the advantages of declarative
programming to imperative programmers. Kaleidoscope initially argued
for a refinement model of constraints, to allow a clean declarative
semantics, and to adjust the imperative language semantics to account
for that (for example, by treating variables as immutable value
streams). Later versions of Kaleidoscope reversed this stance and used
the perturbation model of constraints to be more compatible with an
imperative programming style.

However, Kaleidoscope and related constraint-imperative languages keep
the declarative and imperative paradigms essentially separate,
connected only through shared variables. The two paradigms provide
parallel methods of encapsulation and abstraction, which cannot be
freely mixed. This impedes re-use of modules, as clients can only use
modules in the paradigm for which the implementers provide the
appropriate abstractions.

In this work, we present a language---Babelsberg---that builds on the
ideas from the Kaleidoscope system and other CIP languages. Babelsberg
unifies the methods of encapsulation and abstraction for both the
declarative and the imperative paradigm by using only object-oriented
method definitions for both declarative and imperative code. Loci/R is
integrated into an existing object-oriented VM and, because re-uses
object-oriented messages to declare constraints, allows programmers to
add constraints to existing object-oriented programs in incremental
steps. The performance and semantics of purely object-oriented code is
unaffected in the new VM.

Thus, our current research focuses on the following topics:

...

\section{Background}
\label{sec:background}

\section{Babelsberg}
\label{sec:babelsberg}

In this section, we present our current research progress. It is a
concept for an extension to existing object-oriented languages...

\section{Conclusions and Next Steps}
\label{sec:concl-next-steps}


\bibliography{qiqqa_jot_export,constraints,sigproc}
\end{document}

%%% Local Variables:
%%% mode: latex
%%% TeX-master: t
%%% End:
